\section{Number theory}

\entry{Modular basics}{$O(\log)$ each}{%
\code{power(a, b, m)} is $a^b \bmod m$; it needs $m < 2^{31}$ or the \code{a*a} inside
overflows --- for bigger moduli cast to \code{i128}.\\[0.3ex]
\code{ext\_gcd(a, b, x, y)} returns $g=\gcd$ and fills $ax+by=g$. \code{inv(a, m)}
gives $a^{-1} \bmod m$ for \emph{any} \code{m} with $\gcd(a,m)=1$, and \code{-1} when
none exists. When \code{m} is prime, \code{power(a, m-2, m)} is the same thing and
shorter.\\[0.3ex]
$a^b \bmod m$ with a \emph{gigantic} $b$: read $b$ as a string, reduce it mod
$\varphi(m)$ --- but only add $\varphi(m)$ back (i.e. use $b \bmod \varphi(m) +
\varphi(m)$) when the true $b \ge \varphi(m)$, otherwise you get the wrong answer for
small $b$. That correction is needed exactly when $\gcd(a,m) \ne 1$.}

%LABEL modbasics
\begin{lstlisting}
ll power(ll a, ll b, ll mod) {         // a^b % mod
    a %= mod; ll res = 1;
    for (; b > 0; b >>= 1, a = a * a % mod)
        if (b & 1) res = res * a % mod;
    return res;
}
ll ext_gcd(ll a, ll b, ll &x, ll &y) {     // a*x + b*y = g
    if (b == 0) { x = 1; y = 0; return a; }
    ll x1, y1;
    ll g = ext_gcd(b, a % b, x1, y1);
    x = y1;
    y = x1 - (a / b) * y1;
    return g;
}
ll inv(ll a, ll m) {                 // -1 if none exists
    ll x, y;
    if (ext_gcd(a, m, x, y) != 1) return -1;
    return ((x % m) + m) % m;
}
\end{lstlisting}

\entry{Binomials mod p}{precompute $O(n)$, query $O(1)$ \quad mem $n$}{%
Call \code{precompute()} once, then \code{C(n, k)}. \code{MOD} must be \textbf{prime}
and larger than \code{MAXN}. Returns 0 outside the triangle, so no guards at the call
site.\\[0.3ex]
Non-prime modulus, or $n$ small: build Pascal's triangle instead
(\code{c[i][j] = c[i-1][j-1] + c[i-1][j]}), which needs no inverses. Huge $n$ with small
prime $p$: Lucas' theorem, multiply \code{C} of the base-$p$ digits.\\[0.3ex]
Stars and bars: ways to write $n$ as an ordered sum of $k$ non-negative parts is
\code{C(n+k-1, k-1)}. Catalan $C_n = $ \code{C(2n,n) * inv(n+1)}.}

%DEP modbasics
%LABEL binom
\begin{lstlisting}
const ll MOD = 1000000007, MAXN = 500001;
ll fact[MAXN], inv_fact[MAXN];

void precompute() {
    fact[0] = 1;
    rep(i, 1, MAXN) fact[i] = fact[i-1] * i % MOD;
    inv_fact[MAXN-1] = power(fact[MAXN-1], MOD - 2, MOD);
    down(i, MAXN - 1, 0)
        inv_fact[i] = inv_fact[i+1] * (i+1) % MOD;
}
ll C(ll n, ll k) {
    if (n < 0 || k < 0 || k > n || n >= MAXN) return 0;
    return fact[n] * inv_fact[k] % MOD
                   * inv_fact[n-k] % MOD;
}
\end{lstlisting}

\entry{FFT / NTT (convolution)}{$O(n \log n)$ \quad mem $n$}{%
\code{conv(a, b)} multiplies two polynomials: \code{c[i] = sum a[j]*b[i-j]}, length
\code{a.size()+b.size()-1}. This is also ``for every pair of elements, bucket their
sum'' --- count pairs summing to $s$, subset-sum counts, string matching with
wildcards.\\[0.3ex]
\textbf{NTT}, exact, mod \code{P = 998244353} (a prime with $2^{23} \mid P-1$, root 3).
So the result must be wanted mod that prime, and the transform length must be a power
of two $\le 2^{23}$. Other friendly primes: $1004535809$ (root 3),
$167772161$ (root 3).\\[0.3ex]
If the true coefficients are needed exactly and exceed \code{P}, run it under two
different primes and combine them with the Chinese remainder theorem via
\code{ext\_gcd} (the primes are far too big for \code{CRT.cpp}'s stepping loop). For real/complex input use a
\code{complex<double>} FFT instead --- but it loses precision above $\sim10^{14}$ in the
output, so prefer NTT whenever the answer is modular.}

%LABEL ntt
\begin{lstlisting}
const ll P = 998244353, G = 3;
ll pw(ll a, ll b, ll m = P) {
    ll r = 1; a %= m;
    for (; b; b >>= 1, a = a * a % m)
        if (b & 1) r = r * a % m;
    return r;
}
void ntt(vi& a, bool inv) {
    int n = (int)a.size();
    for (int i = 1, j = 0; i < n; i++) {
        int b = n >> 1;
        for (; j & b; b >>= 1) j ^= b;
        j ^= b;
        if (i < j) swap(a[i], a[j]);
    }
    for (int len = 1; len < n; len <<= 1) {
        ll w = pw(G, (P - 1) / (2 * len));
        if (inv) w = pw(w, P - 2);
        for (int i = 0; i < n; i += 2 * len) {
            ll cur = 1;
            rep(j, 0, len) {
                ll u = a[i+j], v = a[i+j+len] * cur % P;
                a[i+j] = (u + v) % P;
                a[i+j+len] = (u - v + P) % P;
                cur = cur * w % P;
            }
        }
    }
    if (inv) {
        ll ni = pw(n, P - 2);
        for (ll& x : a) x = x * ni % P;
    }
}
vi conv(vi a, vi b) {
    if (a.empty() || b.empty()) return {};
    int need = (int)(a.size() + b.size() - 1), n = 1;
    while (n < need) n <<= 1;
    a.resize(n); b.resize(n);
    ntt(a, false); ntt(b, false);
    rep(i, 0, n) a[i] = a[i] * b[i] % P;
    ntt(a, true);
    a.resize(need);
    return a;
}
\end{lstlisting}

\entry{AND / OR / XOR convolution}{$O(n \log n)$, $n = 2^k$}{%
\code{c[z]} = sum of \code{a[x]*b[y]} over all pairs with \code{x OP y == z}, where OP
is AND (\code{mode 0}), OR (\code{1}) or XOR (\code{2}). The same job as \code{conv} for
$+$, but on bit masks. Pad both arrays with zeros to the same power of two.\\[0.3ex]
Uses: ``number of pairs whose XOR is $z$'' for all $z$ at once; combining two
distributions over subsets; Nim-like games (XOR of piles) summed over choices.
\code{FST} on its own: OR forward = sum over subsets (same as \code{sos}), AND forward
= sum over supersets. Values grow like $n \cdot \max^2$; to work mod a prime, add
\code{\% md} to every line and replace the final \code{/ n} by multiplying with
\code{power(n, md-2, md)}. b5xu's \code{bitwise-xor-convolution}, fixed (the original
does not compile) and extended with KACTL's AND/OR.}

%LABEL fst
\begin{lstlisting}
void FST(vi& a, bool inv, int mode) {    // 0 AND 1 OR 2 XOR
    int n = (int)a.size();
    for (int s = 1; s < n; s *= 2)
        for (int i = 0; i < n; i += 2 * s) rep(j, i, i + s) {
            ll &u = a[j], &v = a[j + s];
            if (mode == 0) u += inv ? -v : v;
            else if (mode == 1) v += inv ? -u : u;
            else tie(u, v) = pii(u + v, u - v);
        }
    if (inv && mode == 2) for (ll& x : a) x /= n;
}
vi bitConv(vi a, vi b, int mode) {     // same size, 2^k
    FST(a, 0, mode); FST(b, 0, mode);
    rep(i, 0, (ll)a.size()) a[i] *= b[i];
    FST(a, 1, mode);
    return a;
}
\end{lstlisting}
